\section{Discussione}
\label{sec:discussion}

\subsection{Guadagno dalla previsione e dalla regolazione}

Il guadagno della previsione si misura confrontando B con A. Sui seed
$50$ e $52$, B raggiunge l'ottimo con budget inferiori tra quelli provati.
Sul seed $51$, invece, peggiora il risultato a budget $50$ e $100$;
sul seed $53$ è migliore ai budget piccoli, ma a $500$ rimane sopra
l'ottimo che A trova. Un MAE inferiore a quello della previsione costante
non implica quindi un miglioramento uniforme dell'ordine di ricerca.

Il guadagno aggiuntivo della regolazione si misura confrontando C con B.
Sul seed $51$, una diversificazione permette di raggiungere l'ottimo
entro $200$ estrazioni, contro il gap del $7{,}01\%$ di B. Sul seed $53$,
due diversificazioni consentono a C di raggiungere l'ottimo entro $500$
estrazioni, mentre B conserva un gap dell'$1{,}55\%$. Sui seed $50$ e
$52$ C non migliora il costo trovato da B ai budget osservati. Un
intervento non equivale quindi a un guadagno: può avvenire quando il giro
ottimo è già stato trovato, senza che la ricerca lo abbia certificato.

La politica effettua deviazioni distanziate anziché sostituire
continuamente l'ordinamento predittivo con il minimo costo parziale.
Il criterio può tuttavia favorire stati poco profondi: il costo già
sostenuto non descrive quello ancora necessario. La stagnazione segnala
un andamento della ricerca, ma non dimostra che il modello stia sbagliando
né identifica quale candidato produrrà un miglioramento.

\subsection{Limiti della valutazione}

I seed $50$ e $51$ hanno contribuito allo sviluppo della politica e
non sono un test indipendente delle scelte del regolatore. I seed $52$
e $53$ sono stati valutati mantenendo fisse tali scelte, ma costituiscono
soltanto due nuovi casi della stessa dimensione e distribuzione dei dati
di addestramento. L'evidenza è preliminare: non dimostra un vantaggio
generale o un comportamento su problemi più grandi. La soglia $100$
non è stata selezionata tramite una ricerca sistematica con un insieme
di validazione dedicato.

Le estrazioni hanno costi diversi tra i metodi: feature, previsioni,
ordinamenti e scansioni della pila richiedono lavoro aggiuntivo. Senza
misure temporali non si può concludere che C sia più veloce. Inoltre,
la scelta del regolatore dipende dal modello, dalla distribuzione dei
dati e dal budget. Il dataset didattico è separato e serve a verificare
il procedimento su esempi leggibili; i suoi errori assoluti non sono
direttamente confrontabili con quelli delle mappe di scala maggiore.

\subsection{Alternative poco promettenti e utilità di ulteriore ricerca}

Una previsione elevata può ritardare un candidato, ma non ne certifica
l'inutilità: la potatura resta basata sul costo accumulato. La
diversificazione può assegnare calcolo a candidati diversi da quelli
che sarebbero estratti normalmente. Il risultato sul seed $53$ mostra
un beneficio del cambiamento di esplorazione rispetto a B.

I dati aggregati non registrano però, per ogni candidato, la previsione
originaria e il miglioramento prodotto dal suo completamento. Non
permettono quindi di identificare sistematicamente stati che apparivano
poco interessanti ma richiedevano ulteriore ricerca: l'ultima domanda
della consegna non è ancora verificata a livello di singolo stato.
Analogamente, il sistema non definisce ancora una distribuzione sulle
alternative né una massa probabilistica irrisolta. Questi aspetti restano
limiti espliciti dell'implementazione, non risultati già misurati.
